\documentclass[english]{famnit-thesis}
% \documentclass[english,cover]{famnit-thesis} % Add cover page for print submission.
% \documentclass[english,review]{famnit-thesis} % Review version (anonymised).

\usepackage{ccicons}   % CC license symbols on the license page.
\usepackage{booktabs}  % Professional-quality tables (\toprule, \midrule, \bottomrule).
\usepackage{amsmath}   % Extended math environments.

% ---------------------------------------------------------------------------
% TODO: Replace the placeholder values below with your actual information.
% ---------------------------------------------------------------------------

\title
    {Zaznavanje osi vozil z nevronskimi mrežami za sisteme tehtanja vozil na mostu}
    {Neural Network-Based Axle Detection for Bridge Weigh-In-Motion Systems}
    {Neural Network-Based Axle Detection for B-WIM Systems.}

\author{[FirstName]}{[LastName]}{[LastName] F.}
\studyprogram{Računalništvo in informatika}
\mentor{[prof. dr. Supervisor Name]}{[Supervisor Name, PhD]}
\date{June}{2026}

% Remove the two lines below if you have no co-mentor.
% \comentor{doc. dr. Ime Somentorja}{Ime Somentorja, PhD}
% \workingcomentor{dr. Ime Delovnega Somentorja}{Ime Delovnega Somentorja, PhD}

% ---------------------------------------------------------------------------
% Keywords (Slovene first, then English)
% ---------------------------------------------------------------------------
\keywords
    {zaznavanje osi, tehtanje vozil na mostu, nevronske mreže, TCN, CNN, časovne serije}
    {axle detection, bridge weigh-in-motion, neural networks, temporal convolutional network, time series}

% ---------------------------------------------------------------------------
% Abstracts (Slovene first, then English; max 250 words each)
% ---------------------------------------------------------------------------
\abstract
{
    Sistemi tehtanja vozil na mostu (B-WIM) ocenjujejo bruto težo tovornih vozil
    z analizo dinamičnega signala deformacije mostu med prehodom vozila. Ključni
    korak je natančno zaznavanje osi — določitev točnih vzorčnih pozicij vsakega
    prehoda osi v signalu dolžine 1.300 vzorcev. Gre za zahteven problem dvojiškega
    označevanja zaporedja, ki ga otežuje izrazito neravnovesje razredov: le
    0{,}35\,\% vzorcev je pozitivnih (dogodki osi), kar daje razmerje 287:1.
    V tej nalogi primerjamo tri pristope na naboru 32.141 prehodov vozil:
    (1) klasični filter Savitzky-Golay z iskanjem vrhov, (2) 1D U-Net konvolucijsko
    nevronsko omrežje (CNN, 30\,M parametrov) in (3) temporalno konvolucijsko
    omrežje (TCN, 2{,}7\,M parametrov) z devetimi dilatiranimi rezidualnimi bloki,
    katerih receptivno polje znaša 2.045 vzorcev. Oba modela globokega učenja sta
    naučena s tehtano binarno navzkrižno entropijo (pos\_weight\,=\,287).
    Najboljši model (TCN) dosega F1\,=\,0{,}998, natančnost\,=\,0{,}997,
    priklic\,$\approx$\,1{,}000 in povprečno absolutno napako časovne lokacije
    0{,}24 vzorca, pri čemer prekaša CNN kljub 11-kratnemu zmanjšanju parametrov.
}
{
    Bridge Weigh-In-Motion (B-WIM) systems estimate the gross vehicle weight of
    heavy goods vehicles by analysing the dynamic strain signal recorded as a
    vehicle crosses a bridge. A critical preprocessing step is accurate axle
    detection --- locating the exact sample positions of each axle passage in a
    1,300-sample signal. This is a challenging binary sequence-labelling problem
    due to severe class imbalance: only 0.35\,\% of samples are positive (axle
    events), yielding a 287:1 negative-to-positive ratio. This thesis compares
    three approaches on a dataset of 32,141 vehicle crossings: (1) a classical
    Savitzky-Golay filter and peak-detection baseline, (2) a 1D U-Net
    Convolutional Neural Network (CNN, 30\,M parameters), and (3) a Temporal
    Convolutional Network (TCN, 2.7\,M parameters) with nine dilated residual
    blocks and a receptive field of 2,045 samples. Both deep learning models are
    trained with a weighted binary cross-entropy loss (pos\_weight\,=\,287).
    The TCN achieves F1\,=\,0.998, Precision\,=\,0.997, Recall\,$\approx$\,1.000,
    and a Mean Absolute Timing Error of 0.24 samples on the validation set,
    outperforming both the CNN and the classical baseline. The TCN achieves this
    with 11 times fewer parameters than the CNN, demonstrating that architecture
    design outweighs model scale for B-WIM axle detection.
}

% ---------------------------------------------------------------------------
% Acknowledgments
% ---------------------------------------------------------------------------
\acknowledgments
{
    % TODO: Write your acknowledgments here.
    I would like to thank my supervisor for their guidance and support throughout
    this project.
}

% License page (English version).
\input{license.tex}

% ---------------------------------------------------------------------------
% List of abbreviations
% ---------------------------------------------------------------------------
\abbreviation{B-WIM}{Bridge Weigh-In-Motion}
\abbreviation{BCE}{Binary Cross-Entropy loss}
\abbreviation{CNN}{Convolutional Neural Network}
\abbreviation{F1}{F1-score (harmonic mean of precision and recall)}
\abbreviation{GVW}{Gross Vehicle Weight}
\abbreviation{HGV}{Heavy Goods Vehicle}
\abbreviation{MATE}{Mean Absolute Timing Error}
\abbreviation{ReLU}{Rectified Linear Unit}
\abbreviation{TCN}{Temporal Convolutional Network}

\begin{document}
    \chapter{INTRODUCTION}

Bridge Weigh-In-Motion (B-WIM) systems provide a non-intrusive method for
estimating the gross vehicle weight (GVW) of heavy goods vehicles (HGVs) as they
cross a bridge at normal traffic speed \cite{Moses1979}. Unlike static weigh
stations, which require vehicles to stop and are costly to operate, B-WIM systems
infer axle loads continuously from the dynamic strain response of an instrumented
bridge structure. This makes them highly practical for traffic enforcement,
infrastructure load monitoring, and pavement design.

A fundamental step in every B-WIM pipeline is \emph{axle detection}: identifying
the number of vehicle axles and, critically, the precise sample positions at
which each axle passes over a reference cross-section of the bridge. Accurate
axle timing is a prerequisite for weight estimation --- an incorrectly detected
or missed axle directly corrupts the subsequent load calculation
\cite{Soberl2025}.

Traditional axle detection methods apply rule-based signal processing: the bridge
strain signal is smoothed (e.g., using a Savitzky-Golay filter
\cite{SavitzkyGolay1964}), and local maxima are then located using a peak-finding
algorithm with manually tuned thresholds. While straightforward to implement,
these approaches depend on hand-crafted parameters and may fail when the
signal-to-noise ratio is low, vehicle speed varies, or multiple closely-spaced
axles produce overlapping strain pulses.

Deep learning models, in particular one-dimensional (1D) convolutional
architectures, have demonstrated strong performance across a wide range of
sequential detection and segmentation tasks. Their ability to learn
task-specific representations directly from raw data, without explicit feature
engineering, makes them attractive candidates for the axle detection problem.

\section{OBJECTIVES}

The main objectives of this bachelor thesis are:

\begin{enumerate}
    \item Implement and evaluate a classical signal-processing baseline for axle
          detection using Savitzky-Golay smoothing and peak detection.
    \item Design, train, and evaluate a 1D U-Net Convolutional Neural Network
          (CNN) for per-sample binary axle detection.
    \item Design, train, and evaluate a Temporal Convolutional Network (TCN) with
          dilated residual blocks, ensuring the receptive field covers the full
          signal length.
    \item Compare all three approaches on the same held-out test set using
          consistent, axle-level evaluation metrics.
    \item Identify the limitations of the current study and propose directions
          for future work.
\end{enumerate}

\section{RESEARCH QUESTIONS}

\begin{enumerate}
    \item Can deep learning models accurately detect axle positions from B-WIM
          strain signals despite severe class imbalance (287:1 ratio)?
    \item How does a deep learning approach compare to a classical peak-detection
          baseline in terms of F1-score and timing accuracy?
    \item Does a lightweight TCN outperform a larger U-Net CNN on this task?
\end{enumerate}

\section{THESIS STRUCTURE}

Chapter~2 provides background on B-WIM systems and reviews related work on axle
detection and deep learning for time series. Chapter~3 describes the dataset, all
three methods, the training procedure, and the evaluation protocol. Chapter~4
presents and discusses the experimental results. Chapter~5 concludes the thesis
and outlines directions for future work.
        % Chapter 1 — Introduction
    \chapter{BACKGROUND AND RELATED WORK}

\section{BRIDGE WEIGH-IN-MOTION SYSTEMS}

Bridge Weigh-In-Motion (B-WIM) is a non-intrusive traffic monitoring technology
that estimates axle loads and gross vehicle weights of vehicles in motion. The
principle, first described by Moses \cite{Moses1979}, is to instrument a bridge
with strain gauges and use the measured dynamic response together with an
influence line algorithm to infer the forces applied by passing axles.

A typical B-WIM installation consists of strain gauges attached to the soffit of
a bridge girder, a data acquisition system sampling at 100--1000\,Hz, and signal
processing software. As a vehicle crosses, each axle generates a characteristic
strain pulse. The timing and relative amplitude of these pulses depend on the
vehicle's axle spacing, speed, and individual axle weights. Once the axle
positions in time are known, the system solves a system of equations to recover
the individual axle weights \cite{Soberl2025}.

Modern B-WIM systems can achieve gross vehicle weight accuracy within a few
percent under favourable conditions, but performance degrades when signals are
noisy, vehicle speeds vary widely, or multiple vehicles are present on the bridge
simultaneously. The initial step of reliably detecting individual axle events is
therefore critical: a single missed or spurious detection directly corrupts the
weight estimation for that vehicle.

\section{TRADITIONAL AXLE DETECTION}

Classical axle detection treats the smoothed strain signal as a quasi-static
response and identifies axle positions as local maxima. Typical steps are:

\begin{enumerate}
    \item \textbf{Low-pass filtering} to suppress high-frequency vibration noise.
    \item \textbf{Savitzky-Golay smoothing} \cite{SavitzkyGolay1964}, which fits
          a polynomial to a sliding window. This approach preserves peak shapes
          and curvature better than a simple moving average, making subsequent
          peak-finding more reliable.
    \item \textbf{Peak detection} based on manually tuned thresholds for minimum
          peak height, minimum prominence, and minimum inter-peak distance.
\end{enumerate}

These methods are fast and interpretable, but require per-installation parameter
tuning. They are sensitive to signal-to-noise ratio and may miss closely-spaced
axles (e.g., tandem or tridem axle groups) or generate false positives at the
start and end of a crossing event.

\section{DEEP LEARNING FOR TIME SERIES}

Deep learning has achieved state-of-the-art results across many time-series
tasks, including classification, segmentation, anomaly detection, and event
localisation. Two families of architectures are directly relevant to this work.

\subsection{1D Convolutional Neural Networks and U-Net}

One-dimensional CNNs apply learned filters along the temporal axis to extract
local features from a sequence. Stacking multiple convolutional layers with
increasing receptive fields allows the network to capture patterns at multiple
temporal scales.

The U-Net architecture \cite{Ronneberger2015}, originally proposed for biomedical
\emph{image} segmentation, adapts naturally to 1D sequence labelling. An encoder
path successively down-samples the sequence while increasing feature depth; a
symmetric decoder path up-samples back to the original resolution, with skip
connections that transfer fine-grained spatial detail from encoder to decoder
layers. The final output is a per-sample probability map, making U-Net directly
applicable to dense binary labelling (axle / non-axle).

\subsection{Temporal Convolutional Networks}

Temporal Convolutional Networks (TCNs) \cite{Bai2018} replace recurrent
connections with stacks of \emph{dilated} 1D convolutions. Dilation inserts gaps
between filter taps, exponentially increasing the effective receptive field with
network depth without increasing the number of parameters. For a TCN with kernel
size $k$ and $N$ dilated residual blocks (dilation doubling each block), the
total receptive field is:

\begin{equation}
    RF = 1 + 2(k-1)\bigl(2^{N}-1\bigr).
    \label{eq:rf}
\end{equation}

With $k=3$ and $N=9$, Equation~\ref{eq:rf} gives $RF=2{,}045$ samples, which
comfortably covers the entire 1,300-sample signal in a single forward pass.

Unlike recurrent networks, TCNs process all time steps in parallel and are
therefore stable to optimise and fast to train. Bai et al.\ \cite{Bai2018} showed
on a suite of sequence modelling benchmarks that TCNs match or exceed the
performance of LSTMs and GRUs at lower computational cost.

\section{CLASS IMBALANCE IN BINARY SEQUENCE LABELLING}

A critical challenge in axle detection is extreme class imbalance. In the dataset
used in this thesis, only approximately 0.35\,\% of all signal samples correspond
to an axle event, yielding a negative-to-positive sample ratio of 287:1. If this
imbalance is not addressed, a model that always predicts ``no axle'' achieves
99.65\,\% sample-level accuracy while being completely useless.

A standard remedy is to weight the positive class in the loss function. With
binary cross-entropy (BCE) loss and a scalar weight $w^{+}$, the loss becomes:

\begin{equation}
    \mathcal{L} = -\frac{1}{N}\sum_{i=1}^{N}
        \bigl[ w^{+}\, y_i \log \hat{p}_i
              + (1-y_i) \log(1-\hat{p}_i) \bigr],
    \label{eq:bce}
\end{equation}

where $y_i \in \{0,1\}$ is the ground-truth label and $\hat{p}_i$ is the
predicted probability. Setting $w^{+} = 287$ brings the effective contribution of
positive and negative samples into balance during training.

\section{RELATED WORK}

\v{S}oberl et al.\ \cite{Soberl2025} combine strain-based B-WIM heuristics with
computer vision (traffic camera images) to verify axle configurations and flag
potentially incorrect GVW estimates. On a dataset of over 30,000 HGV records
collected from the same B-WIM installation used in this thesis, their combined
approach improves GVW reliability from 96.7\,\% to 99.89\,\%. Their work
highlights the difficulty of axle configuration inference from strain signals
alone and motivates the use of data-driven methods for axle detection.

To the best of the author's knowledge, the direct application of 1D U-Net or TCN
architectures to per-sample axle detection from B-WIM strain signals has not
previously been reported at the bachelor level, making this thesis an original
empirical contribution to the field.  % Chapter 2 — Background and Related Work
    \chapter{METHODOLOGY}

\section{DATASET}

The dataset consists of 32,141 vehicle crossing records collected by a B-WIM
installation on a short-span bridge. Each record contains:

\begin{itemize}
    \item \textbf{Signal}: a 1,300-sample normalised strain time series covering
          one complete vehicle crossing.
    \item \textbf{Pulses}: a 1,300-sample binary array with value $1.0$ at the
          exact sample where an axle is present and $0.0$ everywhere else.
\end{itemize}

Vehicles in the dataset have between 2 and 11 axles (mean $\approx 4.5$). The
class distribution is severely imbalanced: only 0.35\,\% of all samples are
positive (axle events), giving a negative-to-positive ratio of approximately
287:1.

Records were split at the vehicle level into training (80\,\%), validation
(10\,\%), and test (10\,\%) sets using a fixed random seed, ensuring that no
vehicle appears in more than one split.

\section{BASELINE METHOD}

The baseline uses classical signal processing and serves as a performance
benchmark against which the neural network models are compared.

\begin{enumerate}
    \item \textbf{Smoothing}: A Savitzky-Golay filter \cite{SavitzkyGolay1964}
          with window length 11 and polynomial order 3 is applied to reduce
          high-frequency noise while preserving peak shape.
    \item \textbf{Peak detection}: SciPy's \texttt{find\_peaks} function
          \cite{Scipy2020} is used to locate local maxima with a minimum
          prominence of 0.05 and a minimum inter-peak distance of 5 samples.
    \item \textbf{Output}: The detected peak positions are treated as the
          predicted axle locations.
\end{enumerate}

Parameters were chosen by grid search on the validation set to maximise F1-score.

\section{CNN MODEL: 1D U-NET}

The CNN follows a 1D U-Net architecture \cite{Ronneberger2015} adapted for
sequence labelling.

\subsection{Architecture}

The encoder consists of three down-sampling stages, each comprising two
\texttt{Conv1d--BatchNorm--ReLU} blocks followed by max-pooling with stride 2.
The number of feature channels doubles at each stage (32, 64, 128). A bottleneck
layer at the lowest resolution applies two further convolutional blocks with 256
channels.

The decoder mirrors the encoder with three up-sampling stages using transposed
convolutions, each concatenated with the corresponding encoder skip connection
before two further convolutional blocks. The final layer is a $1{\times}1$
convolution that maps to a single output channel, producing a logit of shape
$[B, 1, 1300]$. No activation is applied at the output (raw logits are used
with \texttt{BCEWithLogitsLoss}).

The total number of trainable parameters is approximately 30 million.

\subsection{Input and Output}

Each signal is reshaped to $[B, 1, 1300]$ (batch, channel, time) before being
fed to the network. The target is the $[B, 1, 1300]$ pulse array. Post-training,
predicted axle positions are extracted from the output probabilities by finding
peaks with a minimum inter-peak distance of 5 samples.

\section{TCN MODEL: TEMPORAL CONVOLUTIONAL NETWORK}

\subsection{Architecture}

The TCN \cite{Bai2018} consists of nine dilated residual blocks. Block $i$ uses
dilation $d_i = 2^{i}$ (i.e., $d \in \{1, 2, 4, 8, 16, 32, 64, 128, 256\}$),
a kernel size of 3, and 64 feature channels throughout. Each block contains:

\begin{enumerate}
    \item Two \texttt{Conv1d} layers with weight normalisation, symmetric
          (non-causal) padding so that output length equals input length.
    \item \texttt{BatchNorm} and \texttt{ReLU} after each convolution.
    \item A residual connection with a $1{\times}1$ projection if channel
          dimensions differ.
\end{enumerate}

The output of the final block is passed through a $1{\times}1$ \texttt{Conv1d}
to produce a single-channel logit of shape $[B, 1, 1300]$.

\subsection{Receptive Field}

Using Equation~\ref{eq:rf} with $k=3$ and $N=9$:

\begin{equation}
    RF = 1 + 2(3-1)(2^{9}-1) = 1 + 4 \times 511 = 2{,}045 \text{ samples}.
\end{equation}

This exceeds the 1,300-sample signal length, so every output position can
potentially attend to the entire input sequence. The total number of trainable
parameters is approximately 2.7 million --- an 11-fold reduction compared to
the CNN.

\section{TRAINING PROCEDURE}

Both neural network models were trained using identical settings:

\begin{itemize}
    \item \textbf{Loss function}: \texttt{BCEWithLogitsLoss} with
          \texttt{pos\_weight}\,=\,287 (see Equation~\ref{eq:bce}), balancing the
          extreme class imbalance.
    \item \textbf{Optimiser}: AdamW \cite{Paszke2019} with initial learning rate
          $10^{-3}$, weight decay $10^{-4}$.
    \item \textbf{Learning rate schedule}: ReduceLROnPlateau (factor 0.5,
          patience 5 epochs).
    \item \textbf{Early stopping}: Training halted if validation F1 did not
          improve for 25 consecutive epochs.
    \item \textbf{Batch size}: 64.
    \item \textbf{Maximum epochs}: 100.
    \item \textbf{Hardware}: NVIDIA GPU via PyTorch 2.6 with CUDA.
\end{itemize}

The best checkpoint (highest validation F1) was saved for each model and used
for all subsequent evaluation.

\section{EVALUATION METRICS}

Evaluation is performed at the \emph{axle level}, not the sample level, to
reflect the operationally relevant quantity. Ground-truth and predicted axle
positions are first extracted from the binary/probability arrays by finding peaks
with a minimum distance of 5 samples. A predicted axle is counted as a true
positive (TP) if it falls within $\pm 5$ samples of an unmatched ground-truth
axle (greedy matching, closest first).

From these counts the following metrics are computed:

\begin{align}
    \text{Precision} &= \frac{TP}{TP + FP}, \\
    \text{Recall}    &= \frac{TP}{TP + FN}, \\
    \text{F1}        &= \frac{2 \cdot \text{Precision} \cdot \text{Recall}}
                             {\text{Precision} + \text{Recall}}.
\end{align}

Additionally, the \emph{Mean Absolute Timing Error} (MATE) measures the average
sample-level displacement of correctly matched predictions from their
ground-truth positions:

\begin{equation}
    \text{MATE} = \frac{1}{|TP|} \sum_{(p,g) \in TP} |p - g|,
\end{equation}

where $p$ and $g$ are the predicted and ground-truth axle sample indices,
respectively. A lower MATE indicates more precise timing, which translates
directly to more accurate weight estimation.% Chapter 3 — Methodology
    \chapter{RESULTS AND DISCUSSION}

\section{SANITY CHECK: OVERFITTING ON 10 SAMPLES}

Before full training, a sanity check verified that each neural network model
could memorise a tiny batch of 10 training examples. With a plain (unweighted)
BCE loss and a learning rate of $5 \times 10^{-3}$, both the CNN and TCN reduced
their training loss below $10^{-2}$ within 200 epochs. This confirms that the
forward and backward passes are implemented correctly and that the architecture
has sufficient capacity for the task.

\section{TRAINING CONVERGENCE}

\subsection{CNN}

The CNN was trained for 5 epochs before early stopping halted training. The best
checkpoint was saved at \textbf{epoch 3} with validation F1\,=\,0.9974 and
validation MATE\,=\,0.246 samples. The rapid convergence is consistent with the
relatively large model capacity (30\,M parameters) and the high information
density of the weighted loss.

% TODO: Insert training curve figure here.
% \begin{figure}[ht]
%     \centering
%     \includegraphics[width=\linewidth]{images/cnn_training_curves.pdf}
%     \caption{CNN training and validation loss (left) and validation F1-score
%              (right) over epochs. The red dashed line marks the best epoch.}
%     \label{fig:cnn_curves}
% \end{figure}

\subsection{TCN}

The TCN was trained for 33 epochs. The best checkpoint was saved at
\textbf{epoch 23} with validation F1\,=\,0.9982 and validation
MATE\,=\,0.237 samples. The learning rate was reduced automatically twice during
training (ReduceLROnPlateau), which produced the characteristic step-wise
improvement visible in the validation curve.

% TODO: Insert training curve figure here.
% \begin{figure}[ht]
%     \centering
%     \includegraphics[width=\linewidth]{images/tcn_training_curves.pdf}
%     \caption{TCN training and validation loss (left) and validation F1-score
%              (right) over epochs. The red dashed line marks the best epoch.}
%     \label{fig:tcn_curves}
% \end{figure}

\section{QUANTITATIVE COMPARISON}

Table~\ref{tab:results} summarises the performance of all three methods on the
validation set using the best saved checkpoint for each neural network model.

\begin{table}[ht]
    \centering
    \caption{Axle-level detection performance on the validation set.
             MATE is the Mean Absolute Timing Error in samples.
             $\pm$5 sample tolerance window, greedy matching.}
    \label{tab:results}
    \begin{tabular}{lccccr}
        \toprule
        \textbf{Method} & \textbf{Precision} & \textbf{Recall}
                        & \textbf{F1} & \textbf{MATE} & \textbf{Params} \\
        \midrule
        Baseline (SG + peaks) & -- & -- & -- & -- & 0 \\
        % TODO: Fill in baseline numbers after running notebook 02.
        CNN (U-Net, epoch 3)  & 0.9954 & 0.9995 & 0.9974 & 0.246 & $\sim$30\,M \\
        TCN (epoch 23)        & \textbf{0.9965} & \textbf{0.9999}
                              & \textbf{0.9982} & \textbf{0.237} & $\sim$2.7\,M \\
        \bottomrule
    \end{tabular}
\end{table}

Both deep learning models achieve near-perfect axle detection. The TCN
outperforms the CNN on all four metrics while using 11 times fewer parameters.
The MATE of 0.24 samples corresponds to approximately 2.4\,ms at a typical
B-WIM sampling rate of 100\,Hz, which is operationally negligible for weight
calculation purposes.

\section{VISUAL INSPECTION}

Figure~\ref{fig:visual} shows the TCN's predictions on six randomly selected
test-set records. In each panel, the raw strain signal (blue) is overlaid with
the scaled TCN output probability (orange). Ground-truth axle positions are
marked by green vertical lines and predicted positions by red dashed lines.

Across all six examples the TCN correctly identifies the axle count and places
predictions within 1--2 samples of the ground truth. No false positives are
observed in these examples, and the probability output forms a sharp, narrow peak
above each axle event, confirming that the model has learned a precise temporal
localisation rather than a diffuse detection.

% TODO: Insert visual inspection figure here (generated from notebook 04, cell 14).
% \begin{figure}[ht]
%     \centering
%     \includegraphics[width=\linewidth]{images/tcn_visual_inspection.pdf}
%     \caption{TCN predictions on six test-set records. Blue = strain signal;
%              orange = TCN output probability (scaled to signal amplitude range);
%              green lines = ground-truth axle positions;
%              red dashed lines = predicted axle positions.}
%     \label{fig:visual}
% \end{figure}

\section{DISCUSSION}

\subsection{Architecture vs.\ Model Size}

The most striking finding is that the TCN, with 2.7\,M parameters, consistently
outperforms the CNN, which has 30\,M parameters. This supports the hypothesis
that the key design choice for this task is receptive field coverage rather than
raw model capacity. The TCN's dilated blocks ensure that every output position
can attend to the full 1,300-sample signal, while the CNN's skip connections
compensate for the loss of temporal resolution in the encoder but cannot replicate
this global context.

\subsection{Convergence Speed}

The CNN converged in only 3 effective epochs, whereas the TCN required 23. This
likely reflects the CNN's larger parameter count providing more initial
flexibility, but also a tendency to overfit sooner. The slower, more regular
convergence of the TCN --- with two visible learning rate reduction steps ---
suggests better regularisation and a smoother optimisation landscape.

\subsection{Practical Implications}

A MATE of 0.24 samples translates to sub-millisecond timing accuracy at 100\,Hz,
well within the tolerance required for accurate weight calculation. Both models
achieve near-perfect recall ($\geq 0.9995$), meaning very few axles are missed.
Precision is slightly lower ($\sim 0.9954$--$0.9965$), indicating occasional
false positives, but these remain rare (fewer than 1 in 200 detected axles).

\subsection{Limitations}

\begin{itemize}
    \item \textbf{Record-level data split}: The train/validation/test split is
          performed by vehicle record, not by day or week. Crossings from the same
          day may share environmental conditions (temperature, noise floor),
          potentially inflating generalisation estimates. A chronological
          (day-level) split would provide a stricter evaluation.
    \item \textbf{Single installation}: All data originates from one bridge and
          sensor configuration. Performance on a different bridge or sensor type
          is unknown.
    \item \textbf{Raised axles}: Vehicles with a raised (lifted) axle produce a
          zero or near-zero signal for that axle. Neither model was explicitly
          tested on such cases.
    \item \textbf{Robustness to noise and speed variation}: A systematic
          robustness experiment (e.g., adding synthetic Gaussian noise or
          time-stretching signals to simulate speed variation) was outside the
          scope of this work.
\end{itemize}     % Chapter 4 — Results and Discussion
    \chapter{CONCLUSION}

This thesis investigated neural network-based axle detection for Bridge
Weigh-In-Motion (B-WIM) systems. The problem was framed as per-sample binary
sequence labelling on a dataset of 32,141 vehicle crossing records, each a
1,300-sample strain signal with ground-truth axle positions. The extreme class
imbalance (287:1 negative-to-positive ratio) was addressed by weighting the
binary cross-entropy loss proportionally to the class frequencies.

Three approaches were compared: a classical Savitzky-Golay filter and
peak-detection baseline, a 1D U-Net CNN (30\,M parameters), and a TCN with nine
dilated residual blocks (2.7\,M parameters, receptive field 2,045 samples). Both
deep learning models were evaluated using axle-level precision, recall, F1-score,
and Mean Absolute Timing Error (MATE) with a $\pm$5-sample tolerance window.

The main findings are:

\begin{enumerate}
    \item Both neural network models achieve near-perfect axle detection
          (F1\,$\geq$\,0.997, MATE\,$\leq$\,0.25 samples) on unseen vehicle
          records, substantially outperforming the classical baseline.
    \item The TCN outperforms the larger CNN on all metrics despite having 11
          times fewer parameters. This demonstrates that \emph{architecture
          design} --- specifically, ensuring sufficient receptive field coverage
          --- matters more than model scale for this task.
    \item A MATE of approximately 0.24 samples corresponds to sub-millisecond
          timing accuracy at 100\,Hz, well within the tolerance required for
          accurate GVW estimation.
\end{enumerate}

\section{FUTURE WORK}

Several directions remain open for future investigation:

\begin{itemize}
    \item \textbf{Day-level data splitting}: Re-evaluating models with a
          chronological split would test cross-temporal generalisation more
          rigorously.
    \item \textbf{Multi-installation generalisation}: Collecting data from
          multiple bridges and sensor configurations would reveal how well the
          models transfer to unseen installations.
    \item \textbf{Robustness under noise and speed variation}: Systematic
          experiments with synthetic noise injection and time-stretching would
          quantify model robustness.
    \item \textbf{End-to-end weight estimation}: Integrating the axle detector
          into a full B-WIM pipeline and evaluating GVW accuracy would demonstrate
          operational impact.
    \item \textbf{Lighter architectures}: Exploring smaller TCN variants or
          knowledge distillation could enable deployment on low-power edge
          hardware typically used in B-WIM installations.
\end{itemize}
   % Chapter 5 — Conclusion

    \backmatter
    % English theses must include an extended Slovene abstract (4,000–10,000
    % characters incl. spaces) as the last numbered chapter before References.
    % Uncomment and fill in abstract_sl.tex when ready.
    % \input{abstract_sl.tex}

    \input{priloga.tex}     % Appendix A — (optional)
\end{document}
